\chapter{Local Registration: Iterative Closest Point}\label{ch:icp}

Registration finds the transformation that aligns two shapes with partially overlapping geometry.
It is needed to merge the scans of an object (\cref{ch:acq}), to compare a manufactured part with
its CAD model and find defects, to align medical images of the same patient, to track objects, and
to build statistical shape models from collections. This chapter studies the \emph{local}
problem: the two shapes are already roughly aligned and we look for the optimal rigid motion. The
standard algorithm is the iterative closest point (ICP) method of \citet{Besl1992} and
\citet{Chen1992}, one of the most widely used algorithms in computer graphics and in industry.
\Cref{ch:global} removes the assumption of a good initial alignment.

\section{Problem statement}

Let $X=\{\bx_i\}$ (the \emph{source}) and $Y=\{\by_j\}$ (the \emph{target}) be two point clouds.
We look for a rotation $\mR\in SO(3)$ and a translation $\bt\in\R^3$ such that the transformed
source lies on the target. A good alignment is one in which corresponding points are close after
the transformation. If the correspondences $\bx_i\leftrightarrow\by_i$ were known, the problem
would be the least-squares fit
\begin{equation}
\min_{\mR,\bt}\ \sum_{i=1}^N\norm{\mR\bx_i+\bt-\by_i}^2 ,
\label{eq:procrustes}
\end{equation}
which has a closed-form solution. The difficulty is that we know neither the correspondences nor
the transformation, and each would be easy to find given the other.

\section{The ICP algorithm}

ICP alternates between the two easy subproblems (\cref{alg:icp}): given the current
transformation, match each source point to its closest target point; given the matches, solve
\eqref{eq:procrustes}.

\begin{algorithm}[tbp]
\caption{Iterative closest point (point-to-point).}\label{alg:icp}
\Input{source $X$, target $Y$, initial $(\mR,\bt)$, tolerance $\varepsilon$}
\Output{rigid transformation $(\mR,\bt)$}
build a kd-tree on $Y$\;
\Repeat{the decrease of $E$ is below $\varepsilon$ or a maximum number of iterations}{
  \ForEach{$\bx_i\in X$}{$\by_{c(i)}\gets$ nearest neighbour of $\mR\bx_i+\bt$ in $Y$\;}
  reject pairs with distance above a threshold (optional)\;
  $(\mR,\bt)\gets\argmin\sum_i\norm{\mR\bx_i+\bt-\by_{c(i)}}^2$ via \eqref{eq:kabsch}\;
  $E\gets\frac1N\sum_i\norm{\mR\bx_i+\bt-\by_{c(i)}}^2$\;
}
\end{algorithm}

\subsection{Closest points}

The closest point of $\bx$ in $Y$, $\argmin_{\by\in Y}\norm{\by-\bx}$, is found by brute force in
$O(M)$ for a target of $M$ points, so matching $N$ source points costs $O(NM)$ per iteration.
Geometrically, $Y$ partitions space into Voronoi cells and the query asks which cell contains
$\bx$. A kd-tree on $Y$ (\cref{ch:repr}) answers it in $O(\log M)$ expected time, giving
$O(N\log M)$ per iteration; approximate nearest neighbours are faster still and accurate enough.

\subsection{The optimal rigid transformation}

\begin{theorem}[Kabsch, Arun, Umeyama]
Let $\mu_x=\frac1N\sum_i\bx_i$, $\mu_y=\frac1N\sum_i\by_i$ and
$\mC=\sum_i(\by_i-\mu_y)(\bx_i-\mu_x)\T$ with singular value decomposition
$\mC=\mU\mSigma\mV\T$. The minimiser of \eqref{eq:procrustes} is
\begin{equation}
\mR=\mU\,\diag\bigl(1,1,\det(\mU\mV\T)\bigr)\,\mV\T,\qquad \bt=\mu_y-\mR\mu_x .
\label{eq:kabsch}
\end{equation}
\end{theorem}

\begin{proof}
\emph{Translation.} For fixed $\mR$, setting the derivative with respect to $\bt$ to zero gives
$\sum_i2(\mR\bx_i+\bt-\by_i)=\bzero$, hence $\bt=\mu_y-\mR\mu_x$. Substituting, and writing
$\bx_i'=\bx_i-\mu_x$, $\by_i'=\by_i-\mu_y$, the problem reduces to
$\min_\mR\sum_i\norm{\mR\bx'_i-\by'_i}^2$.

\emph{Expanding the error.} Since $\mR\T\mR=\mI$,
$\norm{\mR\bx'_i-\by'_i}^2={\bx'_i}\T\bx'_i-2{\by'_i}\T\mR\bx'_i+{\by'_i}\T\by'_i$. Only the middle term
depends on $\mR$, so we must maximise $\sum_i{\by'_i}\T\mR\bx'_i$.

\emph{The covariance matrix.} Each term is a scalar, equal to its own trace, and with
$\tr(\mA\mB)=\tr(\mB\mA)$: $\sum_i{\by'_i}\T\mR\bx_i'=\tr\bigl(\mR\sum_i\bx'_i{\by'_i}\T\bigr)=\tr(\mR\mC\T)$.

\emph{Maximising the trace.} With $\mC=\mU\mSigma\mV\T$,
$\tr(\mR\mC\T)=\tr(\mR\mV\mSigma\mU\T)=\tr(\mSigma\,\mU\T\mR\mV)$. The matrix
$\mathbf N=\mU\T\mR\mV$ is orthogonal, so its entries satisfy $\abs{N_{kk}}\le1$ and
$\tr(\mSigma\mathbf N)=\sum_k\sigma_kN_{kk}\le\sum_k\sigma_k$. The bound is attained for
$\mathbf N=\mI$, that is $\mR=\mU\mV\T$.

\emph{Avoiding reflections.} If $\det(\mU\mV\T)=-1$, then $\mU\mV\T$ is a reflection. Among
orthogonal $\mathbf N$ with $\det(\mU\mathbf N\mV\T)=+1$ the maximum of $\tr(\mSigma\mathbf N)$ is
$\sigma_1+\sigma_2-\sigma_3$, attained at $\mathbf N=\diag(1,1,-1)$, which gives
\eqref{eq:kabsch}.
\end{proof}

\begin{example}{A planar triangle}{kabsch}
Let $X=\{(0,0,0),(1,0,0),(0,1,0)\}$ and let $Y$ be $X$ rotated by $90^\circ$ about $z$ and
translated by $(1,2,0)$: $Y=\{(1,2,0),(1,3,0),(0,2,0)\}$. Then $\mu_x=(\tfrac13,\tfrac13,0)$,
$\mu_y=(\tfrac23,\tfrac73,0)$ and
\[
\mC=\frac13\begin{pmatrix}1&-2&0\\2&-1&0\\0&0&0\end{pmatrix}=\mR_z(90^\circ)\sum_i\bx'_i{\bx'_i}\T .
\]
The third singular value is zero because the points are coplanar: both the rotation and its
composition with the reflection across the plane fit the data perfectly. The determinant
correction selects the rotation, $\mR=\mR_z(90^\circ)$, and $\bt=\mu_y-\mR\mu_x=(\tfrac23,\tfrac73,0)-(-\tfrac13,\tfrac13,0)=(1,2,0)$.
With fewer than three non-collinear correspondences the rotation is not determined.
\end{example}

\subsection{Convergence}

Let $E=\sum_id(\mR\bx_i+\bt,Y)^2$. The matching step can only decrease each term (each point gets
its closest partner) and the alignment step minimises the sum for the current matches, so $E$
decreases monotonically and, being bounded below, converges. The limit, however, is only a
\emph{local} minimum: ICP converges to the correct alignment only when started close enough to it.
Point-to-point ICP also converges slowly (linearly) when the shapes must slide along each other,
because closest-point pairs pull points towards each other even when they lie on the same flat
region.

\fig{ch07_icp_iterations.png}{Point-to-point ICP aligning two partially overlapping scans of the
bunny, starting from a $25^\circ$ rotation and a translation. Pairs farther than three times the
median distance are rejected, which handles the non-overlapping parts. The rotation error drops from
$25.0^\circ$ to $0.9^\circ$ after 60 iterations.}{fig:icp-iter}

\section{Variants}

\Citet{Rusinkiewicz2001} decompose ICP into stages, each with several options:
\begin{enumerate}
  \item \textbf{Selection} of source points: all points, uniform or random subsampling, or
  \emph{normal-space sampling}, which picks points so that the normals are spread uniformly and
  small features that constrain sliding are not drowned out by large flat regions.
  \item \textbf{Matching}: closest point, closest point along the normal (normal shooting), or
  projection into the range image of the target (very fast for organised scans).
  \item \textbf{Weighting}: uniform, decreasing with distance, or with the compatibility of normals
  $\bn_x\cdot\bn_y$.
  \item \textbf{Rejection} of outliers: pairs farther than a threshold (a multiple of the median or
  of the standard deviation of the distances), pairs with incompatible normals, pairs involving
  boundary points of the target (whose closest-point matches are wrong in partial overlaps), or a
  fixed fraction of the worst pairs (trimmed ICP).
  \item \textbf{Error metric}: point-to-point or point-to-plane (below).
  \item \textbf{Minimisation}: closed form, linearisation, or non-linear least squares.
\end{enumerate}

\subsection{Point-to-plane ICP}

Uneven sampling is a fundamental problem of the point-to-point metric: samples of the source
rarely land exactly on samples of the target, so the residual does not vanish even at the correct
alignment and it penalises motions that slide the surfaces along each other. \Citet{Chen1992}
minimise instead the distance to the tangent plane of the target,
\begin{equation}
E_{\text{plane}}=\sum_i\bigl((\mR\bx_i+\bt-\by_i)\cdot\bn_i\bigr)^2 ,
\end{equation}
where $\bn_i$ is the normal of the target at $\by_i$. Tangential motions are free, so flat regions do
not resist the alignment. There is no closed form, but for small rotations
$\mR\approx\mI+[\br]_\times$ with $\br=(\alpha,\beta,\gamma)$, and
$(\mR\bx_i+\bt-\by_i)\cdot\bn_i\approx\br\cdot(\bx_i\times\bn_i)+\bt\cdot\bn_i-(\by_i-\bx_i)\cdot\bn_i$.
Each pair gives one linear equation in the six unknowns,
\begin{equation}
\mathbf a_i\T\begin{pmatrix}\br\\\bt\end{pmatrix}=b_i,\qquad
\mathbf a_i=\begin{pmatrix}\bx_i\times\bn_i\\\bn_i\end{pmatrix},\qquad b_i=(\by_i-\bx_i)\cdot\bn_i ,
\end{equation}
and the least-squares solution follows from the $6\times6$ normal equations
$\bigl(\sum_i\mathbf a_i\mathbf a_i\T\bigr)\binom{\br}{\bt}=\sum_i\mathbf a_ib_i$ \citep{Low2004}.
The rotation is rebuilt from $\br$ (with Rodrigues' formula, or by projecting $\mI+[\br]_\times$ onto
$SO(3)$) and the step is repeated with new correspondences. Each iteration is a Gauss--Newton
step on the squared distance to the surface, so convergence near the solution is much faster than
for point-to-point (\cref{fig:icp-conv}). The system is singular when the geometry does not
constrain all six degrees of freedom: a plane allows in-plane sliding, a sphere allows any rotation
about its centre, a cylinder allows sliding along and rotating about its axis. The eigenvectors of
$\sum_i\mathbf a_i\mathbf a_i\T$ with small eigenvalues reveal these unstable directions
\citep{Gelfand2005}.

\fig[0.85\textwidth]{ch07_convergence.pdf}{Convergence of ICP on two fully overlapping samplings of
the bunny with a $20^\circ$ initial rotation. Left: root mean square distance to the ground-truth
positions. Right: rotation error. Point-to-plane reaches the noise floor in a handful of
iterations; point-to-point converges linearly.}{fig:icp-conv}

\subsection{Basin of convergence}

Because ICP is a local method, its success depends on the initial misalignment.
\Cref{fig:icp-basin} repeats the registration from random initial rotations of increasing angle:
both variants always succeed up to $45^\circ$; at $60^\circ$ point-to-point succeeds in $88\%$ of
the trials and point-to-plane in all of them; beyond $90^\circ$ point-to-point never succeeds and
point-to-plane only occasionally ($25\%$ at $120^\circ$). Larger misalignments need a global
method (\cref{ch:global}) to provide the initial guess.

\fign{ch07_basin.pdf}{Fraction of successful registrations (final rotation error below $2^\circ$,
eight random trials per angle) against the initial rotation angle.}{fig:icp-basin}

\subsection{Robust and non-rigid extensions}

Least squares is sensitive to outliers and to partial overlap, since every wrong match pulls with
a force proportional to its distance. Robust variants replace the squared distance by a robust
penalty $\rho(d)$ (Huber, Tukey, or $\ell_p$ with $p<1$ as in Sparse ICP \citep{Bouaziz2013}),
which are minimised with iteratively reweighted least squares or ADMM. When the shapes deform
(a moving person, organs), the rigid motion is replaced by a deformation field regularised for
smoothness or local rigidity, and correspondences are optimised jointly with it \citep{Li2008}.

\section{Multi-view registration}

A complete model needs many scans. Registering them sequentially accumulates drift, so after
pairwise alignments a global optimisation distributes the error over all pairs (pose-graph
optimisation), and the aligned scans are merged by reconstruction (\cref{ch:recon}) or by
volumetric fusion of signed distances \citep{Curless1996}.

\begin{keypoints}
\begin{itemize}[leftmargin=*]
  \item ICP alternates closest-point matching (kd-tree, $O(N\log M)$) and the closed-form rigid fit
  $\mR=\mU\diag(1,1,\det\mU\mV\T)\mV\T$, $\bt=\mu_y-\mR\mu_x$.
  \item The error decreases monotonically but converges only to a local minimum; a good initial
  alignment is required.
  \item Point-to-plane ICP linearises the rotation, solves a $6\times6$ system per iteration and
  converges much faster, because surfaces may slide along each other.
  \item Sampling, weighting and outlier rejection make ICP robust to partial overlap and noise.
\end{itemize}
\end{keypoints}

\section*{Exercises}
\begin{exercise}
Show that the optimal translation for a fixed rotation aligns the centroids, and that the problem
then decouples into finding $\mR$ from the centred point sets.
\end{exercise}
\begin{exercise}
Derive the linear system of point-to-plane ICP and explain why it is singular for a plane, a
sphere and a cylinder. Which directions are unconstrained in each case?
\end{exercise}
\begin{exercise}
Extend the closed-form solution to similarity transformations $s\mR\bx+\bt$. Show that the optimal
scale is $s=\tr(\mSigma\mD)/\sum_i\norm{\bx'_i}^2$, with $\mD=\diag(1,1,\det\mU\mV\T)$
\citep{Umeyama1991}.
\end{exercise}
